\subsection*{Consecutive ratios at all nonnegative activities}
\begin{theorem}\label{thm:all-weighted-ratios}
For \(u,v\ge0\), let \(\rho(u,v)\) be the first positive solution of
\(f_{u,v}(\rho)=1/\mu\), before any positive pole. Then
\[
       \lim_{n\to\infty}\frac{Z_{n+1}(u,v)}{Z_n(u,v)}
                 =\frac{1}{\rho(u,v)} .
\]
In particular the ratio for avoidance of both signs tends to
\((\mu-1+\sqrt{\mu^2-6\mu+1})/2\).
\end{theorem}
\begin{proof}
First consider the root limit. Theorem~\ref{thm:joint} and
Corollary~\ref{cor:double} give it in the interior and at the origin.
On the axis \(u=0,\ v>0\), the upper splice used in the proof of
Theorem~\ref{thm:joint} preserves positive avoidance. It can add one
negative occurrence and loses none. Thus
\(Z_{r+s}(0,v)\ge\min(1,v)Z_r(0,v)Z_s(0,v)\).
Reflection gives the other axis. Fekete's lemma and the radius
calculation in that theorem give the root limit \(g=1/\rho(u,v)\)
throughout. Since \(Z_n(u,v)\ge A_n\), this limit exceeds one and
all sufficiently large terms are positive.

Use the four crossing moves of Theorem~\ref{thm:idle-shifts},
from \href{https://arxiv.org/pdf/cond-mat/0008178#page=7}{Jensen's
transfer} and the word description of Albert and Paterson \cite{AP}.
In a signature, list the exposed strand ends across the cut, with a
divider indicating the line. The letter \(U\) inserts the two ends of a
new strand at the divider. The letter \(D\) joins the two ends next to
the divider, provided they belong to distinct strands. The letter \(R\)
moves the last left end to the right and \(L\) moves the first right
end to the left. Pair the internal letters of a meander word starting
immediately after its initial \(U\), omitting its final \(D\).
Every tested site is a triple \(pab\), where \(ab\) is one of these
pairs and \(p\) is the preceding letter. One sign occurs precisely when
\[
         p\in\{U,L\},\quad a=R,\quad b\in\{D,L\},
\]
and the other when \(p\in\{U,R\},\ a=L,\ b\in\{D,R\}\).
Indeed the two short arches require an opening and a closing in
consecutive positions, first on one side and then on the other.
Identify the left side of the list with the upper half-plane of
the arch diagram. Then \(R\) closes above and opens below, so the
first test has weight \(u\), the second \(v\). Reflection exchanges
the tests.

Record the previous letter together with the signature. Let \(W\)
count allowed pairs on this state space, multiplying each pair by
the activity of its sign, or by one when neither sign occurs. At a
pair boundary the number of ends on each side is odd. The initial
state is a single strand across the divider with previous letter
\(U\). It cannot occur after a pair, since any subsequent \(U\)
inserts a new strand in a nonempty signature. Let \(\alpha\) be the
weighted row after the first pair from this initial state, and let
\(\beta\) indicate the single-strand signature across the divider
with any attainable previous letter. Jensen's final closure then gives
\[
                  Z_{N+2}(u,v)=\alpha W^N\beta \qquad (N\ge0).
\]
Here \(W\) acts on the states reachable after at least one pair,
with all original strand pairings retained. All products are
entrywise finite by the four-letter rule.

For every such state we exhibit two returns that preserve its previous
letter and create no marked site. In the table, a bar divides the
exposed ends. The letters \(a,b\) represent ends of pre-existing
strands. Equal digits are the two ends of one local strand. The names
of local strands can be reassigned after each pair. The last column
records the exposed ends after each successive pair of letters.
\begin{center}
\small
\begin{tabular}{@{}llll@{}}
\toprule
Previous & Initial ends & Pairs & Ends after the pairs\\
\midrule
\(\mathrm R\)&\(\mathrm{a|b}\)&
\(\mathrm{UL,DR}\)&\(\mathrm{a11|b,\ a|b}\)\\
& &\(\mathrm{UL,LR,DR}\)&\(\mathrm{a11|b,\ a11|b,\ a|b}\)\\[2pt]
\(\mathrm D\)&\(\mathrm{a|b}\)&
\(\mathrm{LU,UR,DD}\)&\(\mathrm{ab1|1,\ ab1|221,\ a|b}\)\\
& &\(\mathrm{UU,LL,DR,RD}\)&\(\mathrm{a12|21b,\ a1221|b,\ ab1|1,\ a|b}\)\\[2pt]
\(\mathrm U\)&\(\mathrm{ab1|1}\)&
\(\mathrm{UR,DD,LU}\)&\(\mathrm{ab1|221,\ a|b,\ ab1|1}\)\\
& &\(\mathrm{UL,DR,RD,LU}\)&\(\mathrm{ab122|1,\ ab1|1,\ a|b,\ ab1|1}\)\\
\bottomrule
\end{tabular}
\end{center}
Reflection gives the two returns for previous letter \(L\) from the
\(R\) rows. When the previous letter is \(U\), the central ends are
paired. There is at least one older strand, and each side has an even
number of other ends. Hence one side has at least two. Display the
two nearest on that side as \(a,b\), reflecting the \(U\) rows if
necessary. Every other exposed end stays outside the calculation.
For previous letter \(R,L\) or \(D\), a single nearest end on each
side suffices. No step joins \(a\) to \(b\), so \(a\) and \(b\) may
also be paired in the full signature. Each
listed word ends in its prescribed previous letter. Checking its pairs
in order against the two triples above gives weight one, including
the first pair at the incoming boundary.

The return lengths in pairs are \(2,3\) or \(3,4\). Repeating a
return of length 2 gives length 4, so \(W^d\ge I\) entrywise for
\(d=3,4\), for every \(u,v\ge0\).
Fix either \(d\), and for each residue class take a sufficiently large
representative \(r\) that \(\alpha W^r\beta>0\). Then
\[
 \alpha W^{r+dj}\beta
   =\sum_{i=0}^j\binom ji\alpha W^r(W^d-I)^i\beta .
\]
The binomial weights on the right are nonnegative, the zeroth is
positive, and the \(j\)-th root of the left side tends to \(g^d>1\).
Lemma~\ref{lem:binomial-smoothing} gives
\(Z_{n+d}/Z_n\to g^d\) on each residue class and hence on their
finite union. Finally
\[
      \frac{Z_{n+1}}{Z_n}
        =\frac{Z_{n+4}/Z_n}{Z_{n+4}/Z_{n+1}}\ \longrightarrow\ g .
\]
\end{proof}
